\documentclass[12pt, letterpaper]{article}
\date{\today}
\usepackage[margin=1in]{geometry}
\usepackage{amsmath}
\usepackage{hyperref}
\usepackage{cancel}
\usepackage{amssymb}
\usepackage{fancyhdr}
\usepackage{pgfplots}
\usepackage{booktabs}
\usepackage{pifont}
\usepackage{amsthm,latexsym,amsfonts,graphicx,epsfig,comment}
\pgfplotsset{compat=1.16}
\usepackage{xcolor}
\usepackage{tikz}
\usetikzlibrary{shapes.geometric}
\usetikzlibrary{arrows.meta,arrows}
\newcommand{\Z}{\mathbb{Z}}
\newcommand{\N}{\mathbb{N}}
\newcommand{\R}{\mathbb{R}}
\newcommand{\Q}{\mathbb{Q}}
\newcommand{\C}{\mathbb{C}}
\newcommand{\F}{\mathbb{F}}

\newcommand{\Po}{\mathcal{P}}
\newcommand{\Pro}{\mathbb{P}}
\author{Alex Valentino}
\title{Numerical Analysis homework}
\pagestyle{fancy}
\renewcommand{\headrulewidth}{0pt}
\renewcommand{\footrulewidth}{0pt}
\fancyhf{}
\rhead{
	Homework 4\\
	Numerical Analysis
}
\lhead{
	Alex Valentino\\
	Collaborators: Irshad Guliyev (iag38), Jay Patwardhan (jap600)
}
\begin{document}
\begin{itemize}
	\item{Question 1} \textit{Compute the LU factorization of the following matrix
	$$
	A = \begin{bmatrix}
		1 & 2 & 7\\
		3 & 5 & -1\\
		6 & 1 & 4
	\end{bmatrix}	
	$$	in the following two ways:
	}
	\begin{enumerate}
		\item \textit{Gaussian elimination}
		\begin{align*}
			E_1 &= \begin{bmatrix}
			1 & 0 & 0\\
			-3 & 1 & 0\\
			0 & 0 & 1
			\end{bmatrix}\\
			E_1A &= \begin{bmatrix}
			1 & 2 & 7\\
			0 & -1 & -22\\
			6 & 1 & 4
			\end{bmatrix}\\
			E_2 &= \begin{bmatrix}
			1 & 0 & 0\\
			0 & 1 & 0\\
			-6 & 0 & 1
			\end{bmatrix}\\
			E_2E_2A &= \begin{bmatrix}
				1 & 2 & 7\\
				0 & -1 & -22\\
				0 & -11 & -38
			\end{bmatrix}\\
			E_3 &= \begin{bmatrix}
				1 & 0 &0\\
				0 & 1 & 0\\
				0 & -11 & 1 
			\end{bmatrix}\\
			E_3 E_2 E_1 A &= \begin{bmatrix}
			1 & 2 & 7\\
			0 & -1 & -22\\
			0 & 0 & 204
			\end{bmatrix}\\
			&= U
		\end{align*}
		\begin{align*}
		L &= (E_3 E_2 E_1)^{-1}\\
		&= E_1^{-1} E_2^{-1} E_3^{-1}\\
		&= E_1^{-1} \begin{bmatrix}
			1 & 0 & 0\\
			0 & 1 & 0\\
			6 & 0 & 1
			\end{bmatrix}\begin{bmatrix}
				1 & 0 &0\\
				0 & 1 & 0\\
				0 & 11 & 1 
			\end{bmatrix}\\
			&= E_1^{-1} \begin{bmatrix}
			1 & 0 & 0\\
			0 & 1 & 0\\
			6 & 11 & 1
			\end{bmatrix}\\
			&= \begin{bmatrix}
			1 & 0 & 0\\
			3 & 1 & 0\\
			0 & 0 & 1
			\end{bmatrix}\begin{bmatrix}
			1 & 0 & 0\\
			0 & 1 & 0\\
			6 & 11 & 1
			\end{bmatrix}\\
			&= \begin{bmatrix}
			1 & 0 & 0\\
			3 & 1 & 0\\
			6 & 11 & 1
			\end{bmatrix}\\
		\end{align*}
		\item \textit{Doolittle's method}
		$$L = \begin{bmatrix}
		1 & 0 &0\\
		l_{21} & 1 & 0\\
		l_{31} & l_{32} & 1		
		\end{bmatrix}		
		U= \begin{bmatrix}
		u_{11} & u_{12} & u_{13}\\
		0 & u_{22} & u_{23}\\
		0 & 0 & u_{33} 
		\end{bmatrix} 
		A = LU		
		$$
		\begin{align*}
			u_{11}  &= 1\\
			3 &= u_{11} l_{21}\\
			3 &= l_{21}\\
			6 &= u_{11} l_{31}\\
			6 &= l_{31}\\
			u_{12} &= 2\\
			5 &= l_{21} u_{12} + u_{22}\\
			5 &= 6 + u_{22}\\
			-1 &= u_{22}\\
			1 &= u_{12} l_{31} + u_{22} l_{32}\\
			1 &= 2 \cdot 6 - l_{32}\\
			-11 &= -l_{32}\\
			11 &= l_{32}\\
			u_{13} &= 7\\
			-1 &= u_{13}l_{21} + u_{23}\\
			-1 &= 7 \cdot + u_{23}\\
			-22 &= u_{23}\\
			4 &= u_{13} l_{31} + u_{23} l_{32} + u_{33}\\
			4 &= 7 \cdot 6 - 22 \cdot ll + u_{33}\\
			204 &= u_{33}			
		\end{align*}
		Therefore 
		$$\begin{bmatrix}
		1 & 0 & 0\\
		3 & 1 & 0\\
		6 & 11 & 1
\end{bmatrix}		 = \begin{bmatrix}
		1 & 0 &0\\
		l_{21} & 1 & 0\\
		l_{31} & l_{32} & 1		
		\end{bmatrix},		
		\begin{bmatrix}
		1 & 2 & 7\\
		0 & -1 & -22\\
		0 & 0 & 204
		\end{bmatrix} = \begin{bmatrix}
		u_{11} & u_{12} & u_{13}\\
		0 & u_{22} & u_{23}\\
		0 & 0 & u_{33} 
		\end{bmatrix} 					
		$$
	\end{enumerate}
	\item{Question 2} \textit{Solve the triadiagonal system by first computing an LU factorization of the coeffcient matrix and then using the LU factorization to solve the following system:}
		\begin{align*}
	2x_1 + x_2 &= 3\\
	-x_1 + 2x_2 + x_3 &= 2\\
	-x_2 + 2x_3 + x_4 &= 2\\
	-x_3 + 2x_4 + x_5 &= 2\\
	-x_4 + 2x_5 &= 1
	\end{align*}\\
	We're trying to find $(x_1,x_2,x_3,x_4,x_5)$.  Note by the tri-diagonal we need to establish the matrix corresponding to 
	this system of equations.  We have the following:
	$$
	\begin{bmatrix}
	2 & 1 & 0 & 0 & 0\\
	-1 & 2 & 1 & 0 & 0\\
	0 & -1 & 2 & 1 & 0\\
	0 & 0 & -1 & 2 & 1\\
	0 & 0 & 0 &  -1 & 2
	\end{bmatrix} = \begin{bmatrix}
	b_1 & c_1 & 0 & 0 & 0\\
	a_2 & b_2 & c_2 & 0 & 0\\
	0 & a_3 & b_3 & c_3 & 0\\
	0 & 0 & a_4 & b_4 & c_4\\
	0 & 0 & 0 & a_5 & b_5 
	\end{bmatrix}$$
	$$
	\begin{bmatrix}
	3\\ 2\\ 2\\ 2\\ 1
	\end{bmatrix}
	=
	\begin{bmatrix}
	d_1\\
	d_2\\
	d_3\\
	d_4\\
	d_5
	\end{bmatrix}
	$$
	Note by the tridiagonal algorthim we have that 
	$$c_1' = \frac{1}{2}, c_2' = \frac{1}{2 + \frac{1}{2}} = \frac{2}{5},
	c_3' = \frac{1}{2 + \frac{2}{5}} = \frac{5}{12}, c_4' = \frac{1}{2 + \frac{5}{12}} = \frac{12}{29}$$
	$$d_1' = \frac{3}{2}, d_2' = \frac{2 + \frac{3}{2}}{2 + \frac{1}{2}} = \frac{7}{5}, d_3' = \frac{2 + \frac{7}{5}}{2 + \frac{2}{5}} = \frac{17}{12}, d_4' = \frac{2 + \frac{17}{12}}{2 + \frac{5}{12}} = \frac{41}{29}, 
	d_5' = \frac{1 + \frac{41}{29}}{2 + \frac{12}{29}} =\frac{70}{70} = 1$$
	$$ x_5 = 1, x_4 = \frac{41}{29} - \frac{12}{29} = 1, x_3 = \frac{17}{12} - \frac{5}{12} = 1, x_2 = 
	\frac{7}{5} - \frac{2}{5} = 1, x_1 = \frac{3}{2} - \frac{1}{2} = 1$$.
	
	Therefore $(x_1,x_2,x_3,x_4,x_5) = (1,1,1,1,1)$
\end{itemize}
\end{document}
